\clearpage
\appendix
\section{The source working package and the equivalence proof}
\label{app:working-equivalence}
\subsection{Source formulation and comparison conventions}
The following is a mathematical restatement, in our notation, of
\cite[Convention 5.1 and Definition 5.2, pp.~29--30]{oai-ecc}; it is
not a verbatim transcription. The original PDF was checked directly.
Its four numbered clauses are preserved here as (S1)--(S4).

\begin{definition}[Source working package]\label{def:source-working}
Fix an unbounded definable eligible class $\mathcal S$, a finite system
shape, and a set-sized finite-support recipe $F$. For each
$\mu\in\mathcal S$, specify a class $\mathcal C_\mu$ of systems with
each vertex a $\K$-model of size $\mu$. Specify arrows and comparisons
$D$ on bounded sorted tuples. Systems have their full code of internal
maps and fixed named data. Arrows are commuting structural embeddings
with strong vertex maps, possibly subject to a further condition.
\begin{enumerate}[label=\textup{(S\arabic*)},leftmargin=*]
\item \textbf{Structural arrows and unions.}
All structured isomorphisms are arrows; arrows compose and are invariant
under isomorphism. If $X\subseteq Y\subseteq Z$ are eligible and both
inclusions into $Z$ are arrows, so is $X\hookrightarrow Y$.
A nonempty increasing chain of inclusion arrows in $\mathcal C_\mu$
whose union has total size $\mu$ has an eligible union and inclusion
arrows. If its stages include by arrows into one eligible object of
level at least $\mu$, so does its union.
\item \textbf{Presentations and mixed downward closure.}
For every well-order $L$ of eligible size $\mu$, $F(L)\in\mathcal C_\mu$.
Every eligible object is isomorphic to such a presentation at its level.
Elements are labels on finite increasing supports; equality, functions,
relations, and internal maps are evaluated uniformly from labels and
finite order/equality patterns. Substitution along order embeddings gives
arrows, including between different eligible levels. If $\nu\leq\mu$
are eligible, every set of at most $\nu$ vertex occurrences in a
$\mu$-object is contained in a $\nu$-object included by an arrow,
retaining the fixed distinguished data.
\item \textbf{Diagram comparisons.}
For each sorted index pattern, $D$ is an equivalence relation on tuples
at the eligible levels where that pattern has length strictly below the
level. It preserves the atomic information of the whole system code,
is compatible with restriction and reindexing, and is invariant under
all arrows. At any one level, equality is exactly equivalence under an
isomorphism of the containing systems. There is a set bound on the
number of diagrams of each arity, independent of level. A joint diagram
extending the diagram of an actual tuple is realizable over that tuple
when the joint length is bounded at the receiving level; witnesses may
come from larger eligible levels. On term tuples, the comparison depends
only on the labels and joint order/equality pattern of supports, also
across levels. The empty tuple has one diagram.
\item \textbf{Exhaustion.}
Compatible matches of increasing bounded sorted lists, exhausting a
source object, assemble to an arrow when the lists include nested
enumerations of small strong submodels of each source vertex, exhausting
that vertex, and each image enumeration is a strong target submodel.
The source writes these vertex submodels as $V_n$.
AEC unions and the common-upper-bound axiom establish component
strongness; any extra scheme-specific arrow requirement must be checked
for this construction when proving the package for that scheme.
\end{enumerate}
The source's main convention uses strong-limit levels of countable
cofinality, apart from explicitly preliminary levels. It does not
explicitly quantify an ordinal indexing the opening lists of (S4).
Our comparison below handles the written countable interpretation, and
also arbitrary \emph{ordinal-indexed nested} lists at countable-cofinality
source levels. It does not replace chains by arbitrary directed families.
\end{definition}

\begin{remark}[Scope of the equivalence claim]\label{rem:equivalence-scope}
For the comparison, use the \emph{same} concrete system code on both
sides for a nonempty finite system shape, with $|\tau^{\mathrm{gen}}|\leq\mu_0$ and
$\mu_0\geq\LS(\K)$. These bounds normalize the usual AEC convention
that its LS number bounds its vocabulary. They are explicit hypotheses
of our comparison, not inequalities stated in source Definition~5.2.
One can restrict to an eligible tail above these fixed bounds, but doing
so changes the minimum and discards preliminary levels; it is not an
equivalence preserving the original class of levels.

For $\operatorname{Diag}_{\mathcal I}(\K)$, assume the source's specified
internal maps are structural embeddings. Alternatively use the general
system category of Definition~\ref{def:working-coding}, where this
restriction is unnecessary. A graph is represented by its free category;
relations between paths are retained if they are part of the chosen
shape. Neither a new commutativity condition nor additional named
skeleton data may be imposed during translation.
\end{remark}

\subsection{Consequences of lifted comparison}
\begin{lemma}[Common targets and basic type properties]\label{lem:la-basic}
Under (A1)--(A4), every two eligible objects have an eligible common
target. Lifted types have atomic information, restriction and reindexing,
and invariance under $H$. The empty type is unique, and every eligible
level is categorical as a class of systems.
\end{lemma}
\begin{proof}
Present $X,Y$ by well-orders $L,M$. Embed both into a well-order obtained
by concatenating them and appending a well-order of sufficiently large
eligible size. The induced arrows, transported through the charts, give
a common target. For the empty tuple these arrows witness $E$ with the
identity automorphism, so its type is unique. The orbit axiom then gives
same-level system categoricity.

In a witness $hf\ba=g\bb$, the maps $hf$ and $g$ are structural
embeddings of the full codes. They preserve and reflect atomic formulas,
including equalities, sorts, internal maps, and named constants. Thus
the tuples have the same atomic information. Restricting or reindexing
the tuple keeps the same witness, as long as the resulting arity is
bounded. For $f\in H(X,Y)$, the maps $f,\operatorname{id}_Y$ and
$\operatorname{id}_Y$ witness
$E((\ba;X),(f\ba;Y))$. Equivalence then gives invariance of all
comparisons under applying arrows. Notice that the direct witness is
automatic, whereas transitivity and the orbit characterization are not.
\end{proof}

\begin{lemma}[Mixed downward closure and transport]\label{lem:la-transport}
Every set $A$ of at most $\nu$ occurrences in $X\in\mathcal C_\mu$,
where $\nu\leq\mu$ are eligible, is contained in the image of an arrow
from a $\mathcal C_\nu$-object. Every eligible $X$ has an arrow into
some object at any prescribed larger eligible level.
\end{lemma}
\begin{proof}
In a well-ordered chart for $X$, choose finite supports for the elements
of $A$ and, if needed, for the fixed named data. There are at most $\nu$
positions, since the code has at most $\mu_0\leq\nu$ names. Extend them
to a suborder $J$ of size $\nu$. Its presentation maps by an arrow into
the chart; transporting and taking its image gives the required
inclusion. A well-order of cardinality $\nu$ has order type in
$[\nu,\nu^+)$, as required by the definition of $\mathcal C_\nu$.
For upward transport append an order of the desired larger cardinality
to the chart for $X$ and use its inclusion arrow.
\end{proof}

\begin{lemma}[Uniform bound, realization, and term naturality]
\label{lem:la-realization}
Lifted types satisfy the bound, realization, and naturality clauses of
(S3).
\end{lemma}
\begin{proof}
Fix a sorted arity of cardinality $\theta$, choose eligible $\kappa>\theta$,
and fix $M\in\mathcal C_\kappa$. Any tuple of that arity can be
transported to level $\kappa$: use the preceding lemma to shrink its
containing object if larger, or to embed it upward if smaller. Arrow
invariance preserves its type, and categoricity takes it into $M$.
Thus there are at most $\kappa^\theta$ types of that sorted arity.
This is a bound on equivalence classes, not a claim that a class of
all realizations is literally a set.

For realization let $\ba\in X\in\mathcal C_\mu$ and
$(\bb,\bc)\in Y$ have total length bounded at both levels, with
$\operatorname{tp}_{\mathrm{la}}(\ba;X)=
\operatorname{tp}_{\mathrm{la}}(\bb;Y)$.
Transport the entire joint tuple into $Y'\in\mathcal C_\mu$ by the
preceding lemma. Its image $(\bb',\bc')$ retains its type.
The orbit axiom supplies $s:Y'\cong X$ with $s\bb'=\ba$.
Then $\bd=s\bc'$ realizes the requested joint type over $\ba$.
The phrase ``possible at a larger level'' in (S3) is included in this
argument; a comparison is only defined where its arity is bounded.

For naturality, identify the joint supports of the two term records
by the specified order isomorphism. Amalgamate their two indexing
well-orders over those supports: in each cut concatenate the two
remaining well-ordered intervals, retaining each common support point.
This is again a well-order. Append positions if necessary to give it
eligible size. The two induced arrows carry equal labels with their
matched supports to literally the same tuple. They witness $E$ with
the identity automorphism. This works across different levels and
includes supports for named data. It uses the \emph{joint} pattern,
not merely separate order types of individual supports.
\end{proof}

\subsection{The generated AEC}\label{app:completion}
\begin{lemma}[Small designated decompositions]\label{prop:eligible-mu0-union}
Under (S1)--(S2), or under the corresponding alternate axioms, the poset
\[
 P_{\mu_0}(X)=\{A\in\mathcal C_{\mu_0}:A\hookrightarrow X\text{ is an arrow}\}
\]
is $\mu_0^+$-directed and has union $X$. Its transition inclusions are
designated arrows.
\end{lemma}
\begin{proof}
Each occurrence belongs to some member by mixed downward closure.
For any family of at most $\mu_0$ members, their total sorted universe
has size at most $\mu_0$. Mixed downward closure gives a small designated
subobject containing this union. Coherence makes all the inclusions
from the original members arrows. This proves $\mu_0^+$-directedness,
not merely pairwise directedness, and also the assertion about the union.
\end{proof}

\begin{theorem}[Standard upward completion, applied to the small slice]
\label{thm:small-completion}
Under (S1)--(S2) and the coding bounds of
Remark~\ref{rem:equivalence-scope}, the AEC in $\mu_0$
$(\mathcal C_{\mu_0},H_{\mu_0})$ has an upward completion
$\K^{\mathrm{gen}}$ with LS number $\mu_0$, and
\[
 \K^{\mathrm{gen}}_{\mu_0}=(\mathcal C_{\mu_0},H_{\mu_0})
\]
on both objects and embeddings. Every eligible object belongs to this
AEC, and every designated arrow is a strong embedding in it. Its concrete
forgetful functor to the ambient system category is faithful and preserves
directed colimits.
\end{theorem}
\begin{proof}
Let $A\leq_0 B$ mean that $A\subseteq B$ and this inclusion belongs to
$H_{\mu_0}$. Isomorphism invariance, reflexivity, transitivity, coherence,
and unions and smoothness at size $\mu_0$ follow from (S1). Thus this is
an AEC in $\mu_0$. Apply the standard upward-completion theorem
\cite[Definition 1.0.16 and Fact 1.0.18]{jarden-shelah-frames}, citing
Shelah's Chapter II, Lemma 1.23. It forms directed unions of the small
models with coherent small-arrow witnesses for the strong relation;
it does not inherit all ambient morphisms. The theorem gives the AEC
axioms, LS number, and exact recovery of the small slice. We invoke this
established completion theorem rather than claim to reprove it here.

The preceding lemma expresses each eligible $X$ as such a directed union,
so $X\in\K^{\mathrm{gen}}$. For an arrow $f:X\to Y$, identify $X$
with $f[X]\subseteq Y$. For every $A\in P_{\mu_0}(X)$ its inclusion
into $Y$ is an arrow by composition. Whenever such $A$ is contained in
$B\in P_{\mu_0}(Y)$, coherence gives $A\leq_0 B$.
Consequently the two small-subobject presentations fit together as the
coherent presentations defining $X\leq_{\mathrm{gen}}Y$: one may use
$P_{\mu_0}(X)$ inside $P_{\mu_0}(Y)$, with the inherited small-arrow order.
This also gives, equivalently, directed pairs of small source and target
pieces covering the whole arrow.

Every small map is componentwise strong. Directed unions and the
common-upper-bound property in $\K$ therefore make all the induced
vertex maps strong in the completion; internal symbols and fixed names
are retained on the sorted unions. Colimits in this AEC and in the
ambient system category are computed on the same carriers. Thus the
forgetful functor is well-defined, concrete, faithful, and preserves
directed colimits. In the diagram version, unions also retain the
internal structural-embedding property.
\end{proof}

\begin{remark}[What an arbitrary ambient completion does not give]
The completion theorem has no conclusion that all generated embeddings
between eligible objects satisfy their original geometric conditions.
The inclusion $H(X,Y)\subseteq\operatorname{Emb}_{\K^{\mathrm{gen}}}(X,Y)$
is all the translation needs. Conversely, any AEC as in (D1) with the
prescribed small slice is recovered from that slice by upward completion:
its small strong submodels cover every model by LS and coherence, and
coherent small-map unions are strong by the directed-union theorem for
AECs. Thus the completion is an auxiliary presentation, not a freely
chosen source of additional eligible arrows.
\end{remark}

\subsection{The two translations}
\begin{lemma}[The source relation is forced]\label{lem:source-la}
For a source working package, $D$-equality is exactly lifted automorphism
comparison $E$ formed from its eligible objects and designated arrows.
In particular, there is no additional choice of cross-level equivalence
once those objects and arrows are fixed and the package axioms hold.
\end{lemma}
\begin{proof}
If $hf\ba=g\bb$ witnesses $E$, invariance under the arrows $f,g,h$
and transitivity of $D$ give $D(\ba;X)=D(\bb;Y)$.
Conversely, choose any common eligible target $Z$ for $X,Y$, with arrows
$f,g$, using their well-ordered presentations. Arrow invariance gives
$D(f\ba;Z)=D(g\bb;Z)$. The same-level orbit clause supplies
$h\in\operatorname{Aut}(Z)$ with $hf\ba=g\bb$, witnessing $E$.
The arity is bounded in $Z$ because it was bounded in the source
objects. Since $D$ is an equivalence relation, so is $E$.
\end{proof}

\begin{lemma}[Short and arbitrary size-preserving fixed-level chains]
\label{lem:long-fixed-chains}
The chain clauses in (A3) imply those in (S1).
\end{lemma}
\begin{proof}
Let $(X_i)_{i<\delta}$ be a nonempty inclusion chain at level $\mu$
whose union has size $\mu$. If it has a final stage, that is the union.
If $\operatorname{cf}(\delta)\leq\mu$, use a cofinal subchain of length
at most $\mu$, apply (A3), and use composition for omitted stages.
If $\operatorname{cf}(\delta)>\mu$, choose a stage containing each
occurrence of the union. Those at most $\mu$ stage indices have a common
bound below $\delta$, so the chain has already stabilized to its union.
The common-upper-bound conclusion follows by the same argument.
The converse follows immediately by taking chains of length $<\mu^+$.
\end{proof}

\begin{lemma}[Countable reduction of nested exhaustion]
\label{lem:countable-exhaustion}
Suppose $\operatorname{cf}(\mu)=\omega$, $|X|=\mu$, and
$(A_\alpha)_{\alpha<\delta}$ is an increasing family of sets of size
$<\mu$ with union $X$. Then $\delta$ has countable cofinality, and a
cofinal countable subchain still exhausts $X$. Any simultaneously nested
vertex resolutions and compatible matches restrict to this subchain.
\end{lemma}
\begin{proof}
There can be no final stage. Put $\kappa=\operatorname{cf}(\delta)$
and pass to a cofinal subchain. If $\kappa>\mu$, the indices of the
first occurrences of all elements have a bound below $\delta$,
contradicting boundedness of a stage. Thus $\kappa\leq\mu$.
If $\mu=\aleph_0$, it follows that $\kappa=\omega$.
Otherwise $\mu$ is singular, so the regular cardinal $\kappa$ is $<\mu$.
The union has size $\mu$, hence the cardinalities of the stages are
unbounded below $\mu$; a bound $\lambda<\mu$ would give union size at
most $\kappa\cdot\lambda<\mu$ (enlarge $\lambda$ to be infinite).
Choose countably many stages whose sizes are cofinal in $\mu$.
If $\kappa>\omega$, their indices have a bound below $\delta$, forcing
one stage to have size at least $\mu$, a contradiction. Hence
$\kappa=\omega$. Cofinal restriction retains all the required
containments, matches, strong images, and unions.
\end{proof}

\begin{theorem}[Equivalence on the specified scope]\label{thm:working-equivalence}
Fix the concrete coding and cardinal bounds of
Remark~\ref{rem:equivalence-scope}. At eligible levels of countable
cofinality, the source package of Definition~\ref{def:source-working}
and the alternate definition of Definition~\ref{def:working-scheme}
are equivalent, preserving the eligible objects, designated arrows,
presentation recipe, and all tuple comparisons. Use the diagram version
when internal maps are structural embeddings, or the general system
version otherwise. Interpret (S4) for countable or, equivalently here,
ordinal-indexed nested exhaustion with the corresponding nested
vertex witnesses.
\end{theorem}
\begin{proof}
\emph{Source to alternate.} Theorem~\ref{thm:small-completion} constructs
$\K^{\mathrm{gen}}$ and embeds all the eligible data in it, recovering
the small slice exactly. Use $F$ as the finite-support system-code
recipe $\Psi$; no all-order AEC blueprint theorem is required. By (S2)
the eligible objects are exactly the isomorphism closures in (D3).
Retain the original $H$. Lemma~\ref{lem:source-la} identifies $D$ with
$E$, proving (A1) and (A4). Presentation arrows give (A2).
Structural closure gives (A3), and (S4) gives (A5) upon taking the
shared enumerations and their partial maps $p\restriction A_n$.

\emph{Alternate to source.} Forget the auxiliary completion, retain
$\mathcal C,H,F$, and define $D$-equality to be equality of lifted types.
The structural assertions of (S1) follow from (A3) and
Lemma~\ref{lem:long-fixed-chains}. The presentation assertions of (S2)
are (D2)--(D3) and (A2); its mixed downward part is
Lemma~\ref{lem:la-transport}. Lemmas~\ref{lem:la-basic} and
\ref{lem:la-realization}, together with (A1) and (A4), prove every
clause of (S3), including the empty type and the uniform bound.

For (S4), compatible matches define well-defined sorted partial
injections because they preserve equality. Their union is total on
$X$. Every finite list of arguments of a function or relation eventually
lies in one stage; atomic agreement there shows preservation and
reflection in the union. Function values also eventually occur, and
names are preserved. Thus the union is a structural embedding of the
whole system code. Its vertex witnesses and their strong images are
exactly those in (A5), so it is in $H$. If the exhaustion is indexed by
an arbitrary ordinal, apply Lemma~\ref{lem:countable-exhaustion} first.
Conversely countable exhaustion is a particular such exhaustion.

The translations leave objects and $H$ unchanged, and
Lemma~\ref{lem:source-la} shows that they recover the original $D$.
The completion is canonical in the sense explained after
Theorem~\ref{thm:small-completion}. This proves equivalence of the
packages, not higher-level fullness of $H$ in the completion.
\end{proof}

\subsection{What remains open or outside this equivalence}
\begin{enumerate}[leftmargin=*]
\item The exact diagram formulation does not cover non-embedding
internal maps without the stated ambient-category change. This is a
scope difference, not an unproved implication that we silently assume.
\item The language and LS bounds at the original minimum are explicit
normalizations. If they fail for a source descriptor, the theorem applies
to a suitable eligible tail, not to all its preliminary levels unchanged.
\item Countable reduction resolves the indexing issue for nested
ordinal chains at countable-cofinality levels. It says nothing about
arbitrary directed matches or preliminary source levels of other
cofinalities; a stronger exhaustion convention there must be added
on both sides to obtain the same comparison.
\item Whether the full package forces all eligible generated strong
embeddings to lie in $H$ remains open in these notes. It is neither
assumed nor needed in the equivalence proof.
\item Formal arbitrary-order evaluations remain distinct from actual
AEC membership. Also, changing the eligible class by later closure
refinements requires the source's separate supply-of-levels hypotheses;
mere unboundedness is not closedness under those refinements.
\end{enumerate}
